\chapter{2022年上海卷(秋)}

\section{填空题}

\begin{problem}
若复数 \(z=1+\i\)(其中 \(\i\) 为虚数单位),则 \(2z=\fillinblank{}\).
\end{problem}

\begin{problem}
双曲线 \(\frac{x^2}{9}-y^2=1\) 的实轴长为\fillinblank{}.
\end{problem}

\begin{problem}
函数 \(y=\cos^2x-\sin^2x+1\) 的最小正周期为\fillinblank{}.
\end{problem}

\begin{problem}
设 \(a\) 为常数,若行列式 \(\begin{vmatrix}a&1\\3&2\end{vmatrix}\) 的值与行列式 \(\begin{vmatrix}a&0\\4&1\end{vmatrix}\) 的值相等,则 \(a=\fillinblank{}\).
\end{problem}

\begin{problem}
若圆柱的高为 \(4\),底面积为 \(9\pi\),则该圆柱的侧面积为\fillinblank{}.
\end{problem}

\begin{problem}
若实数 \(x\),\(y\) 满足 \(x-y\le0\),\(x+y-1\ge0\),则 \(z=x+2y\) 的最小值为\fillinblank{}.
\end{problem}

\begin{problem}
在 \((3+x)^n\) 的二项展开式中,若 \(x^2\) 的系数是常数项的 \(5\) 倍,则 \(n=\fillinblank{}\).
\end{problem}

\begin{problem}
已知 \(f(x)\) 为奇函数,且 \(f(x)=\begin{cases}
a^2x-1, & x<0,\\
x+a, & x>0,
\end{cases}\) 则实数 \(a=\fillinblank{}\).
\end{problem}

\begin{problem}
为了检测学生的身体素质指标,需从游泳类 \(1\) 项,球类 \(3\) 项,田径类 \(4\) 项,共 \(8\) 项项目中随机抽取 \(4\) 项进行测试,每一类都被抽到的概率为\fillinblank{}(结果用最简分数表示).
\end{problem}

\begin{problem}
已知等差数列 \(\{a_n\}\) 的公差不为零,\(S_n\) 为其前 \(n\) 项和. 若 \(S_5=0\),则 \(S_1\),\(S_2\),\ldots,\(S_{100}\) 这 \(100\) 个数中所有不同数值的个数为\fillinblank{}.
\end{problem}

\begin{problem}
已知平面非零向量 \(\bs{a}\),\(\bs{b}\),\(\bs{c}\) 的模均为 \(\lambda\),若 \(\bs{a}\cdot\bs{b}=0\),\(\bs{a}\cdot\bs{c}=2\),\(\bs{b}\cdot\bs{c}=1\),则 \(\lambda=\fillinblank{}\).
\end{problem}

\begin{problem}
设函数 \(y=f(x)\),定义域为 \([0,+\infty)\),值域为 \(A\),且对定义域中任意实数 \(x\) 都成立 \(f(x)=f\left(\frac1{x+1}\right)\). 若 \(\{y\mid y=f(x),x\in[0,a]\}=A\),则实数 \(a\) 的取值范围是\fillinblank{}.
\end{problem}

\section{单选题}

\begin{problem}
若集合 \(A=[-1,2)\),\(B=\Z\),则 \(A\cap B=\)(\quad)

\choices
    {\(\{-2,-1,0,1\}\)}
    {\(\{-1,0,1\}\)}
    {\(\{-1,0\}\)}
    {\(\{-1\}\)}
\end{problem}

\begin{problem}
若实数 \(a\),\(b\) 满足 \(a>b>0\),则下列不等式中,恒成立的是(\quad)

\choices
    {\(a+b>2\sqrt{ab}\)}
    {\(a+b<2\sqrt{ab}\)}
    {\(\frac a2+2b>2\sqrt{ab}\)}
    {\(\frac a2+2b<2\sqrt{ab}\)}
\end{problem}

\begin{problem}
如图,在正方体 \(ABCD-A_1B_1C_1D_1\) 中,\(P\),\(Q\),\(R\),\(S\) 分别为棱 \(AB\),\(BC\),\(BB_1\),\(CD\) 的中点,连接 \(A_1S\),\(B_1D\). 对于空间任意两点 \(M\),\(N\),若线段 \(MN\) 上不存在也在线段 \(A_1S\),\(B_1D\) 上的点,则称 \(M\),\(N\) 两点``可视'',则下列选项中与点 \(D_1\)``可视''的点为(\quad)

\bitmapfigure[max width=0.42\linewidth]{img/2022/shanghai/q15_fig1.png}

\choices
    {点 \(P\)}
    {点 \(B\)}
    {点 \(R\)}
    {点 \(Q\)}
\end{problem}

\begin{problem}
设 \(\Omega=\{(x,y)\mid (x-k)^2+(y-k^2)^2=4|k|,\ k\in\Z\}\),有结论:
\begin{circlelist}
\item 存在直线 \(l\),使得 \(\Omega\) 中不存在点在 \(l\) 上,但存在点在 \(l\) 两侧;
\item 存在直线 \(l\),使得 \(\Omega\) 中存在无数个点在 \(l\) 上.
\end{circlelist}
关于以上两个结论,正确的判断是(\quad)

\choices
    {\(\circled{1}\)成立,\(\circled{2}\)成立}
    {\(\circled{1}\)成立,\(\circled{2}\)不成立}
    {\(\circled{1}\)不成立,\(\circled{2}\)成立}
    {\(\circled{1}\)不成立,\(\circled{2}\)不成立}
\end{problem}

\section{解答题}

\begin{problem}
如图,已知三棱锥 \(P-ABC\),底面 \(ABC\) 为等边三角形,\(O\) 为 \(AC\) 的中点,\(AP=AC=2\),且 \(PO\perp\) 底面 \(ABC\).
\begin{enumerate}
\item 求三棱锥 \(P-ABC\) 的体积 \(V\);
\item 若 \(M\) 为 \(BC\) 中点,求 \(PM\) 与平面 \(PAC\) 所成角的大小.
\end{enumerate}

\bitmapfigure[max width=0.34\linewidth]{img/2022/shanghai/q17_fig1.png}
\end{problem}

\begin{problem}
设 \(a\) 为常数,\(f(x)=\log_3(a+x)+\log_3(6-x)\).
\begin{enumerate}
\item 若将函数 \(y=f(x)\) 的图像向下平移 \(m\)(\(m>0\))个单位后,其图像经过 \((3,0)\),\((5,0)\) 两点,求实数 \(a\) 和 \(m\) 的值;
\item 若 \(a>-3\) 且 \(a\ne0\),解不等式 \(f(x)\le f(6-x)\).
\end{enumerate}
\end{problem}

\begin{problem}
如图,\(AD=BC=6\),\(AB=20\),\(\angle DAB=\angle ABC=120^\circ\),\(O\) 为 \(AB\) 中点,曲线 \(CD\) 上任意一点到点 \(O\) 的距离均相等,\(P\),\(M\),\(Q\) 为曲线 \(CD\) 上的点,且 \(MO\perp AB\),点 \(P\) 与点 \(Q\) 关于直线 \(OM\) 对称.
\begin{enumerate}
\item 若点 \(P\) 与点 \(C\) 重合,求 \(\angle POB\) 的大小;
\item 点 \(P\) 在何位置时,五边形 \(CPMQD\) 的面积 \(S\) 取到最大值,并求出该最大值.
\end{enumerate}

\bitmapfigure[max width=0.58\linewidth]{img/2022/shanghai/q19_fig1.png}
\end{problem}

\begin{problem}
设椭圆 \(\Gamma:\frac{x^2}{a^2}+\frac{y^2}{b^2}=1\)(\(a>b>0\)),\(F_1(-\sqrt2,0)\),\(F_2(\sqrt2,0)\) 为 \(\Gamma\) 的焦点,\(A\) 为 \(\Gamma\) 的下顶点,\(M\) 为直线 \(l:x+y-4\sqrt2=0\) 上一点.
\begin{enumerate}
\item 若 \(a=2\),且 \(AM\) 的中点在 \(x\) 轴上,求点 \(M\) 的坐标;
\item 若直线 \(l\) 与 \(y\) 轴的交点为 \(B\),直线 \(AM\) 经过点 \(F_2\),且在 \(\triangle ABM\) 中有一内角的余弦值为 \(\frac35\),求 \(b\) 的值;
\item 若 \(\Gamma\) 上一点 \(P\) 到 \(l\) 的距离为 \(d\),且 \(|PF_1|+|PF_2|+d=6\),求 \(d\) 的最小值.
\end{enumerate}
\end{problem}

\begin{problem}
在数列 \(\{a_n\}\) 中,\(a_1=1\),\(a_2=3\),对任意 \(n\in\N^*\) 且 \(n\ge2\),均存在正整数 \(i\in[1,n-1]\),满足 \(a_{n+1}=2a_n-a_i\).
\begin{enumerate}
\item 求 \(a_4\) 所有可能的值;
\item 命题 \(p\):若 \(a_1\),\(a_2\),\ldots,\(a_8\) 成等差数列,则 \(a_9<30\),证明 \(p\) 是真命题,同时写出 \(p\) 的逆命题 \(q\),并判断命题 \(q\) 是真命题还是假命题,说明理由;
\item 若 \(a_{2m}=3^m\)(\(m\in\N^*\)),求数列 \(\{a_n\}\) 的通项公式.
\end{enumerate}
\end{problem}

